\section{Related Work}
\label{sec:related_work}

\paragraph{Reinforcement learning with verifiable rewards.}
RLHF~\citep{ouyang2022training} aligns language models with human
preferences through learned reward models trained on preference
annotations. DeepSeek-R1~\citep{deepseekr1} and
DeepSeekMath~\citep{shao2024deepseekmath} replaced learned preference
models with programmatic correctness checks, establishing the
reinforcement learning with verifiable rewards (RLVR) paradigm for
mathematical reasoning. GRPO~\citep{shao2024deepseekmath} computes
group-relative advantages without a separate critic network, reducing
memory cost relative to PPO while maintaining competitive performance on
math benchmarks. These methods supervise only the final answer and treat
the intermediate reasoning trace as opaque. TopoPRM retains the
verifiable reward paradigm but extends it to process-level supervision:
the topology and continuity rewards are computed by deterministic
programs, yet they provide step-level structural feedback that
outcome-only RLVR cannot offer.

\paragraph{Process reward models.}
\citet{lightman2023lets} introduced process reward models (PRMs) that
assign scalar rewards to individual reasoning steps, using human
annotations of step-level correctness as training signal.
Math-Shepherd~\citep{wang2024math} reduces annotation cost by estimating
step-level value through Monte Carlo rollouts from each intermediate
state. Recent work has moved toward \emph{generative} PRMs:
ThinkPRM~\citep{thinkprm2025} uses long chain-of-thought verification
to judge each step, achieving strong out-of-domain performance with
orders of magnitude fewer process labels;
GenPRM~\citep{genprm2026} incorporates code verification into the
verification chain, enabling a 1.5B verifier to outperform GPT-4o on
ProcessBench.
RLKD~\citep{rlkd2025} introduces a Generative Structure Reward Model
that captures multi-branch reasoning structure for distillation.
These approaches all rely on learned or LLM-based verifiers.
TopoPRM takes a complementary approach: rather than
scoring individual steps with a neural model, it extracts structural
properties of the full reasoning graph (acyclicity, connectivity,
directional consistency) using deterministic rules, requiring no
annotations, no LLM inference, and no additional model parameters.

\paragraph{Graph-structured reasoning.}
Graph-of-Thought~\citep{besta2024graph} organizes LLM reasoning into
graph structures through prompt engineering, enabling exploration of
multiple reasoning paths during inference.
GRP~\citep{grp2025} trains models to produce graph-structured outputs
with explicit cognitive labels, requiring specialized output formats and
task-specific decoding procedures. TopoPRM is closest to this line and
explicitly builds on its process-aware perspective, but differs in two
practical choices: (1) it keeps standard text generation format and
extracts implicit DAGs post-hoc, and (2) it emphasizes compatibility with
stock RLVR training stacks through plugin-level verifiable rewards. This
design avoids graph-annotation requirements and architecture changes while
retaining structure-aware optimization signals.

\paragraph{Reward hacking and reasoning compression.}
Reward hacking, where models exploit auxiliary reward signals without
genuine task improvement, is a persistent problem in RL-based
training~\citep{Denison2024SycophancyTS}. In the GRPO setting,
DAPO~\citep{dapo2026} introduces dynamic sampling to stabilize training,
and difficulty-aware GRPO~\citep{difficultyGRPO2026} adjusts
optimization pressure based on problem difficulty. On the compression
side, OPSDC~\citep{opsdc2026} uses on-policy self-distillation with
concise privileged teachers to compress verbose reasoning by 40--57\%
while improving accuracy;
reverse-KL distillation~\citep{hinton2015distilling} reduces verbose
chain-of-thought outputs into shorter traces. TopoPRM
addresses reward hacking through Stratified
Clipping Advantage Estimation (SCAE), which hierarchically prioritizes answer
correctness over auxiliary process rewards. Our TVSD compression stage
shares the on-policy self-distillation principle with OPSDC but uses
\emph{deterministic topology diagnostics} rather than instruction-conditioned
teacher logits as the densification signal.

\paragraph{On-policy distillation for reasoning.}
On-policy distillation (OPD) trains a student on its own generations
while a teacher provides dense
supervision~\citep{minillm2023,agarwal2024gkd}.
OPD has become a standard post-training primitive in production systems
including Qwen3~\citep{qwen3_2025}, MiMo~\citep{mimo2026},
and DeepSeek-V4~\citep{deepseekv4_2026}.
ExOPD~\citep{exopd2026} casts OPD as dense KL-constrained RL with
reward extrapolation, enabling students to surpass their teachers.
SRPO~\citep{srpo2026} unifies GRPO and self-distillation via
sample routing.
TopoPRM connects to this line through its TVSD compression stage,
which is an on-policy distillation procedure where the densification
signal comes from deterministic DAG diagnostics rather than teacher
logits or learned reward models. This makes it compatible with any
base model without cross-tokenizer or cross-architecture constraints.
